% Part: computability
% Chapter: recursive-functions
% Section: sequences

\documentclass[../../../include/open-logic-section]{subfiles}

\begin{document}

\olfileid[hi]{cmp}{rec}{seq}
\olsection{अनुक्रम}

आदिम पुनरावर्ती फलनों का समुच्चय उल्लेखनीय रूप से सुदृढ़ है। किंतु
\emph{अनुक्रमों} को सँभालने का पर्याप्त साधन विकसित कर लेने पर हम और भी
अधिक कर पाएँगे। हम प्राकृतिक संख्याओं के परिमित अनुक्रमों की प्राकृतिक
संख्याओं से निम्न प्रकार पहचान करेंगे: अनुक्रम
$\langle a_0, a_1, a_2, \dots, a_k \rangle$ उस संख्या के संगत है:
\[
p_0^{a_0+1} \cdot p_1^{a_1+1} \cdot p_2^{a_2+1} \cdot \dots \cdot
p_k^{a_k+1}.
\]
हम घातांकों में एक जोड़ते हैं, ताकि यह सुनिश्चित हो कि, मसलन, अनुक्रमों
$\langle 2, 7, 3\rangle$ और $\langle 2, 7, 3, 0, 0 \rangle$ के
संख्यात्मक कूट भिन्न हों। हम $0$ और~$1$, दोनों को रिक्त अनुक्रम का कूट
मान सकते हैं; निश्चितता के लिए $\emptyseq$ से~$0$ निरूपित करें।

अनुक्रमों का यह कूटन इसलिए काम करता है कि अंकगणित का तथाकथित मूल प्रमेय
मान्य है: प्रत्येक प्राकृतिक संख्या $n \ge 2$ को ठीक एक ढंग से निम्न रूप
में लिखा जा सकता है:
\[
n = p_0^{a_0} \cdot p_1^{a_1} \cdot \dots \cdot p_k^{a_k}
\]
जहाँ $a_k \ge 1$। यह सुनिश्चित करता है कि प्रतिचित्रण $\tuple{}(a_0,
\dots, a_k) = \tuple{a_0, \dots, a_k}$ एकैकी है: भिन्न अनुक्रम भिन्न
संख्याओं पर भेजे जाते हैं; प्रत्येक संख्या के संगत अधिकतम एक अनुक्रम है।

अब हम दिखाएँगे कि किसी अनुक्रम की लंबाई निर्धारित करने, उसका $i$वाँ अवयव
निर्धारित करने, किसी अनुक्रम के अंत में एक अवयव जोड़ने और दो अनुक्रमों का
संलग्नन करने की संक्रियाएँ सभी आदिम पुनरावर्ती हैं।

\begin{prop}
  फलन $\len{s}$, जो अनुक्रम $s$ की लंबाई लौटाता है, आदिम पुनरावर्ती है।
\end{prop}

\begin{proof}
  $R(i, s)$ को निम्न प्रकार परिभाषित संबंध मानें:
  \[
  R(i, s) \text{ तभी और केवल तभी जब }
  p_i \mid s \land p_{i+1} \nmid s.
  \]
  $R$ स्पष्टतः आदिम पुनरावर्ती है। जब भी $s$ किसी अरिक्त अनुक्रम का कूट
  हो, अर्थात्
  \[
  s = p_0^{a_0+1} \cdot \dots \cdot p_{k}^{a_{k}+1},
  \]
  तब $R(i,s)$ मान्य है यदि $p_i$ ऐसा सबसे बड़ा अभाज्य है कि
  $p_i \mid s$, अर्थात् $i = k$। इस प्रकार $s$ की लंबाई $i+1$ तभी और
  केवल तभी है जब $p_i$,~$s$ को विभाजित करने वाला सबसे बड़ा अभाज्य हो;
  इसलिए हम रख सकते हैं:
  \[
  \len{s} =
  \begin{cases}
    0 & \text{यदि $s = 0$ अथवा $s = 1$} \\
    1 + \bmin{i < s}{R(i, s)} & \text{अन्यथा}
  \end{cases}
  \]
  हम यहाँ परिबद्ध न्यूनीकरण का उपयोग कर सकते हैं, क्योंकि केवल एक $i$
  ऐसा है जो $R(i,s)$ को संतुष्ट करता है जब $s$ किसी अनुक्रम का कूट हो,
  और यदि $i$ अस्तित्व में है तो वह स्वयं~$s$ से छोटा है।
\end{proof}

\begin{prop}
  फलन $\fn{append}(s,a)$, जो $a$ को अनुक्रम $s$ के अंत में जोड़ने का परिणाम
  लौटाता है, आदिम पुनरावर्ती है।
\end{prop}

\begin{proof}
  $\fn{append}$ को इस प्रकार परिभाषित किया जा सकता है:
  \[
  \fn{append}(s,a) =
  \begin{cases}
    2^{a+1} & \text{यदि $s = 0$ अथवा $s = 1$} \\
    s \cdot p_{\len{s}}^{a+1} & \text{अन्यथा।}
  \end{cases}
  \]
\end{proof}

\begin{prop}
  फलन $\fn{element}(s,i)$, जो $i$ क्रमांक का अवयव $s$ में से लौटाता है (जहाँ
  आरंभिक अवयव को $0$वाँ कहते हैं), अथवा $0$ लौटाता है यदि $i$, $s$ की
  लंबाई से बड़ा या उसके बराबर हो, आदिम पुनरावर्ती है।
\end{prop}

\begin{proof}
ध्यान दें कि $a$, $i$ क्रमांक पर~$s$ का अवयव तभी और केवल तभी है जब
$p_i^{a+1}$,~$p_i$ की वह महत्तम घात हो जो~$s$ को विभाजित करती है, अर्थात्
$p_i^{a+1} \mid s$ किंतु $p_i^{a+2} \nmid s$। अतः:
  \[
  \fn{element}(s,i) =
  \begin{cases}
    0 & \mbox{यदि $i \geq \len{s}$} \\
    \bmin{a < s}{(p_i^{a+2} \nmid s)} & \text{अन्यथा।}
  \end{cases}
  \]
\end{proof}

ऊपर परिभाषित फलनों के औपचारिक नामों का उपयोग करने के बदले हम अधिक
संक्षिप्त संकेत-पद्धति प्रस्तुत करते हैं। हम $(s)_i$ लिखेंगे, न कि
$\fn{element}(s,i)$, और निम्न को संक्षेप में लिखने के लिए
$\tuple{s_0, \dots, s_k}$ का उपयोग करेंगे:
\[
\fn{append}(\fn{append}(\dots \fn{append}(\emptyseq,s_0)
\dots),s_k).
\]
ध्यान दें कि यदि $s$ की लंबाई~$k$ है, तो $s$ के अवयव
$(s)_0$, \dots,~$(s)_{k-1}$ हैं।

\begin{prop}
दो अनुक्रमों का संलग्नन करने वाला फलन $\fn{concat}(s,t)$ आदिम
पुनरावर्ती है।
\end{prop}

\begin{proof}
  हमें ऐसा फलन $\fn{concat}$ चाहिए जिसमें यह गुण हो:
  \[
    \fn{concat}(\tuple{a_0, \dots, a_k}, \tuple{b_0, \dots, b_l}) =
    \tuple{a_0, \dots, a_k, b_0, \dots, b_l}.
  \]
  हम एक ``सहायक'' फलन $\fn{hconcat}(s,t,n)$ का उपयोग करेंगे, जो प्रथम $n$
  अवयव~$t$ से लेकर~$s$ के साथ संलग्न करता है। इस फलन को आदिम
  पुनरावर्तन द्वारा इस प्रकार परिभाषित किया जा सकता है:
  \begin{align*}
    \fn{hconcat}(s,t,0) & = s\\
    \fn{hconcat}(s,t,n+1) & = \fn{append}(\fn{hconcat}(s,t,n),(t)_n)
    \intertext{तब हम $\fn{concat}$ को इस प्रकार परिभाषित कर सकते हैं:}
    \fn{concat}(s,t) & = \fn{hconcat}(s,t,\len{t}).
  \end{align*}
\end{proof}

हम $s \concat t$ लिखेंगे, न कि $\fn{concat}(s,t)$।

किसी अनुक्रम के संख्यात्मक कूट को उसकी लंबाई और उसके महत्तम अवयव के पदों
में परिबद्ध कर पाना हमारे लिए उपयोगी होगा। मान लें कि $s$, लंबाई~$k$ का
ऐसा अनुक्रम है जिसका प्रत्येक अवयव किसी संख्या~$x$ से छोटा या उसके बराबर
है। तब~$s$ के अधिकतम $k$ अभाज्य गुणनखंड हैं, प्रत्येक अधिकतम~$p_{k-1}$
है, और प्रत्येक की घात अधिकतम $x+1$ है,~$s$ के अभाज्य गुणनखंडन में। दूसरे
शब्दों में, यदि हम परिभाषित करें:
\[
\fn{sequenceBound}(x,k) = p_{k-1}^{k \cdot (x+1)},
\]
तो ऊपर वर्णित अनुक्रम~$s$ का संख्यात्मक कूट अधिकतम
~$\fn{sequenceBound}(x,k)$ है।

अनुक्रमों पर ऐसा परिबंध होने से हमें परिबद्ध खोज का उपयोग करके नए फलन
परिभाषित करने का ढंग मिलता है। मसलन, हम परिबद्ध खोज का उपयोग करके
$\fn{concat}$ परिभाषित कर सकते हैं। हमें केवल उस वस्तु (संलग्न अनुक्रम की
संख्या) का आदिम पुनरावर्ती \emph{विनिर्देश} और यह परिबंध लिखना है कि खोज
कहाँ तक करनी है। निम्न काम करता है:
\begin{align*}
  \fn{concat}(s,t) = {} & \bmin{v < \fn{sequenceBound}(s+t,\len{s} +
    \len{t})}{} \\
  & \quad(\len{v} = \len{s} + \len{t} \land {}\\
    & \qquad \bforall{i < \len{s}}{((v)_i = (s)_i) \land {} \\
      & \qquad \bforall{j < \len{t}}{((v)_{\len{s}+j} = (t)_j)})}
\end{align*}

\begin{prob}
दिखाएँ कि ऐसा आदिम पुनरावर्ती फलन~$\fn{sconcat}(s)$ है जिसमें यह गुण है:
\[
\fn{sconcat}(\tuple{s_0, \dots, s_k}) = s_0 \concat \dots \concat s_k.
\]
\end{prob}

\begin{prob}
दिखाएँ कि ऐसा आदिम पुनरावर्ती फलन~$\fn{tail}(s)$ है जिसमें यह गुण है:
\begin{align*}
  \fn{tail}(\emptyseq) & = 0 \text{ और}\\
  \fn{tail}(\tuple{s_0, \dots, s_{k}}) & = \tuple{s_1, \dots, s_{k}}.
\end{align*}
\end{prob}

\begin{prop}
  \ollabel{prop:subseq}
  फलन $\fn{subseq}(s, i, n)$, जो $s$ का लंबाई~$n$ का वह उपानुक्रम
  लौटाता है जो $i$वें अवयव से आरंभ होता है, आदिम पुनरावर्ती है।
\end{prop}

\begin{proof}
  अभ्यास।
\end{proof}

\begin{prob}
  \olref[cmp][rec][seq]{prop:subseq} सिद्ध करें।
\end{prob}
  
\end{document}
